\begin{table}[t]
\centering\scriptsize
\setlength{\tabcolsep}{3.5pt}
\renewcommand{\arraystretch}{0.92}
\caption{Main results: selective escalation on the internal test set
and five public speech-QA benchmarks (delivered answer accuracy, \%,
native full-duplex sessions, one live arm per tier). The probe is
fixed; tier thresholds are label-free calibration quantiles per
language (the internal tiers realize 8/25/52\% escalation;
realized rates per pool are in Appendix~\ref{app:transferdetail}).
Matched-rate random mixes the always-local and measured
always-escalate arms at each pool's realized rate. avg and $\Delta$
(average gain over always-local) cover five external pools in the
live block and four in the NVDA block. The relay-free bound scores the
expert's returned text directly. The
bottom block evaluates our latest three-layer
NemotronLabs-VoiceChat-11B probe by re-mixing local and cached expert
outcomes after replay through its native frame protocol
(Appendix~\ref{app:nvda}). External columns retain cached local outcomes
and the prior 2,481-row fit. The Internal column judges the same pass-3
answers on all 240 queries and uses a 2,258-row refit that excludes
these test IDs (with cached pass-2 training labels). All policies leave
the 17 queries without an answer-onset read local, so the 15/30/50\%
scoreable-query budgets realize 14/28/47\% overall. This Internal result
is post-selection because probe architecture selection included this
split. AlpacaEval, reported separately on its
1--5 scale, changes from 3.68 to 4.01 at 44\% escalation on MiniCPM
(matched random 3.91; always-escalate 4.20 against 4.95 for the
expert text, a spoken-length cap); the pass-3 NVDA ranking changes
from 3.55 to 4.51 versus 4.23 for random.}
\label{tab:transfer}
\begin{tabular}{lcccccccc}
\toprule
& Internal & TriviaQA & WebQ & Llama Q. & SD-QA & Reason.\ zh & avg & $\Delta$ \\
\midrule
\multicolumn{9}{l}{\emph{native full-duplex live sessions (delivered accuracy)}} \\
always-local          & 40.8 & 62.8 & 52.8 & 82.4 & 48.0 & 51.0 & 59.4 & --- \\
\;+ gate, conservative & 46.7 & 60.4 & 59.2 & 78.8 & 54.0 & 62.4 & 63.0 & +3.6 \\
\;+ gate, balanced     & 54.2 & 70.0 & 62.4 & 82.4 & 65.0 & 66.3 & 69.2 & +9.8 \\
\;\textbf{+ gate, aggressive} & \textbf{62.9} & \textbf{88.4} & \textbf{75.2} & \textbf{85.6} & \textbf{82.5} & \textbf{72.8} & \textbf{80.9} & \textbf{+21.5} \\
\;matched random, aggressive & 56.1 & 81.5 & 68.9 & 84.1 & 73.7 & 67.5 & 75.2 & +15.8 \\
always-escalate (measured) & 70.4 & 96.0 & 79.6 & 91.6 & 87.0 & 83.7 & 87.6 & +28.2 \\
expert text as returned (relay-free bound) & 76.5 & 96.4 & 82.3 & 93.2 & 88.5 & 86.1 & 89.3 & +29.9 \\
\midrule
\multicolumn{9}{l}{\emph{NVDA three-layer answer-onset probe (pass-3 native-frame replay)}} \\
always-local          & 19.6 & 35.8 & 37.6 & 70.8 & 31.0 & --- & 43.8 & --- \\
\;+ gate @15\%        & 31.7 & 48.4 & 48.8 & 80.7 & 42.0 & --- & 55.0 & +11.2 \\
\;+ gate @30\%        & 44.6 & 61.0 & 59.2 & 85.8 & 54.0 & --- & 65.0 & +21.2 \\
\;+ gate @50\%        & 60.4 & 76.8 & 69.6 & 89.7 & 66.5 & --- & 75.7 & +31.9 \\
\;matched random @50\% & 53.5 & 65.7 & 60.2 & 82.0 & 60.0 & --- & 67.0 & +23.2 \\
always-expert (ceiling) & 91.7 & 95.5 & 82.8 & 93.1 & 89.0 & --- & 90.1 & +46.3 \\
\bottomrule
\end{tabular}
\end{table}


\begin{figure}[!t]
  \centering
  \includegraphics[width=\linewidth]{revision_accuracy_latency}
  \caption{\textbf{Answer-content accuracy and separately measured serving time.}
  Left: accuracy against expert call rate. Internal uses post-hoc
  weights of 0.25 for knowledge/math and 1 otherwise; all 240 queries
  remain included. Table~\ref{tab:transfer} retains the original weights;
  external curves are unweighted. The dashed line is the random
  endpoint mixture; the horizontal line is the measured always accuracy.
  Right: mean and P50 of signed server PCM-ready TTFA from a separate
  serving run on the same query IDs, with original query weights and
  a different frozen gate snapshot (Appendix~\ref{app:ttfa}). Statistics
  use completed sessions; negative times retain early responses.
  The panels do not establish joint accuracy--latency performance.
  L/C/B/A/W denote local/conservative/balanced/aggressive/always.
  Table~\ref{tab:latency} gives counts and tails for all five pools.}
  \label{fig:dualview}
\end{figure}

\subsection{Internal accuracy and timing cost}
\label{sec:live}
\label{sec:liveloop}
\label{sec:livecurve}
The upper MiniCPM block of Table~\ref{tab:transfer} reports
answer-content accuracy on the fixed 240-query internal test set:
40.8\% locally, then 46.7\%, 54.2\%, and 62.9\% for conservative,
balanced, and aggressive. Their realized escalation rates are
8.3\%, 25.0\%, and 51.7\% (Appendix~\ref{app:transferdetail}).
Aggressive exceeds the table's random reference of 56.1\% by
6.8 percentage points. Per-tier internal comparisons are in
Table~\ref{tab:live}. Aggressive recovers 74.6\% of the measured
local-to-always accuracy gap while using 51.7\% expert calls.
This measures selection value relative to the recorded endpoints.

Gains vary across the fixed internal strata: factual QA improves by
52.5 points, while math accuracy is unchanged. Of 60 knowledge
questions, 43 exceed the expert's input window. These descriptive
results identify an input constraint and task differences, without
attributing the accuracy gap to truncation or changing the test
mixture (Appendix~\ref{app:internal-strata}).

The separate serving run measures when the first answer audio becomes
available (Figure~\ref{fig:dualview}). Internal TTFA P50 is 0.70\,s
locally, 0.89\,s at conservative, 1.61\,s at balanced, 2.87\,s at
aggressive, and 6.67\,s at always. Aggressive has mean 4.59\,s and
P95 18.14\,s; always has mean 8.01\,s and P95 25.86\,s.
Local mean is $-0.61$\,s because 17 of 239 completed sessions begin
answering before the input ends. Such early responses do not establish
successful turn-taking. Across 5,950 attempts, 5,904 complete and
46 fail; distributions condition on completion and retain outliers
(Table~\ref{tab:latency}). These measurements characterize server
response time, while correctness under this causal input protocol
and client-audible delay remain unmeasured.
